大语言模型能力受参数规模、训练数据量、数据质量及领域配比共同影响；扩模、提质和延长上下文均受算力约束。本文沿“数据刻画—规律建模—资源优化—能力预测”的路径建立四问模型，并区分来源观测、半合成与外推估算的数据用途。

针对问题一，将附件 A 的 22 个质量字段展开为 25 个方向统一的指标，在 A1 拟合 CRITIC 权重及 Huber 稳健中心，构造文档主质量分并取领域中位数。A1 与 A2/A3 对应领域的分布未检出显著差异（KS 检验 $p=0.313/0.220$）。评分后按单指标与六来源分层诊断；在 A1/A2/A3，冻结筛查规则给出的高分歧候选比例为 $5.00\%/4.73\%/4.99\%$，多来源对立候选比例为 $1.36\%/1.58\%/1.77\%$。这些比例不是质量错误率，诊断不回写主分。来源一致性加权的实验对照分与主分的秩相关为 $0.9547/0.9207/0.9395$，不替代主口径。

通过 A16/A18 将七个质量领域桥接至 17 个配比领域，并用中心对数比变换估计配比与 13 个验证域交叉熵的关系。线性主模型在 1M 留出集上的平均逐域 $R^2=0.772$；经过训练配比信赖域约束，由该模型选出的等权候选配比，其加权质量为 $0.6624$，相对均匀配比的预测降损为 $9.47\%$，尚未经重新训练验证。GBDT 作为非线性对照，1M 留出的平均逐域 $R^2=0.923$。冻结 1M 线性模型后，60M 和 1B 的目标排序 Spearman 分别为 $0.467$、$0.677$；用 60M 标定尺度修正后，在独立 1B 检验上的平均逐域 MAE 从 $2.953$ 降至约 $0.77$，但绝对损失仍存在明显偏差。A12--A15 的估算表不作为真实大模型验证。

针对问题二，以 B1 校准经典标度律，在 B6 的 360 个半合成点拟合质量项；指定的幂次主式在 B7 新增 90 点的留出 RMSE 为 $0.053$，线性对照为 $0.042$，A、B 质量口径的锚点仍属建模假设。问题三通过解析消去训练数据量并数值求解预算配置；问题四分层桥接损失与能力，比较规模扩张和非规模技术进步。问题三、四的既有数值尚依赖旧接口，须贯通重算；缺少人工质量真值与真实大规模训练复验时，稳健性和代理模型预测收益不能解释为实际性能提升。

\keywords{数据质量评价\quad 领域配比\quad 广义标度律\quad 算力约束优化}
